\documentclass[a4paper,11pt]{article}

\usepackage{revy}
\usepackage[utf8]{inputenc}
\usepackage[T1]{fontenc}
\usepackage[danish]{babel}


\revyname{MatematikRevyen}
\revyyear{2025}
% HUSK AT OPDATERE VERSIONSNUMMER
\version{0.1}
\eta{$2$ minutter}
\status{Færdig}

\title{En Sketchy Sketch}
\author{Thais '22, Augusta '20}

\begin{document}
\maketitle

\begin{roles}
\role{I}[Snow] Instruktør
\role{S}[Conrad] Studerende 
\role{D}[Vitus] Drugdealer 
\end{roles}

\begin{props}
\item{Stol}
\end{props}


\begin{sketch}

\scene D står i mellemgangen med en lang frakke på og ser lusket ud. S kommer ind af døren og går hen til D. Igennem hele sketchen snakkes der ikke men hviskes højlydt.  

\says{S}[Hvisker] Psst, Hey D. Jeg hører du sælger.

\says{D} Hvor har du hørt det fra?

\says{S} Min ven Rasmus. Han siger det er dig der har de gode sager.

\says{D} Ah Rasmus. Ja han er en mine faste kunder.

\says{S} Så du sælger?

\says{D} Ja, den er god nok. Hvad er det helt præcis du søger?

\says{S} Du ved, eksamen nærmer sig og jeg er på mit aller sidste forsøg. Jeg har brug for noget der kan... hjælpe mig lidt på vej, forstår du? 

\says{D} Jeg forstår. Og jeg tror lige jeg har det der kan hjælpe.

\scene D rækker ind under jakken så man får en ide om at hun/han har ting og sager i inderlommen.

\says{D} Men jeg må advare dig, det er ikke for sarte sjæle. 

\says{S} Jeg er parat til at prøve alt. Bare det virker.

\says{D} Om det virker. Dette produkt er af den højeste kvalitet.

\says{S} Okay fedt!

\says{D} Det vil gøre dine linjer skarpere, dine tanker klarere. Bruger du det rigtigt, kan det bære dig gennem selv de sværeste beviser. Det er ikke magi, men det føles sådan.

\says{S} Wow! 

\says{D} Det er ikke for amatører. Det er derfor alle de store professorer bruger det.

\says{S} Virkelig? Er det japansk? 

\says{D} Det er japansk. 

\says{S} Okay jeg er solgt, men hvordan ved jeg, det er den ægte vare du står med?

\says{D} Du mærker det når du bruger det. Mærk på posen, det er den fineste hvide pulver! 

\says{S} Altså det er Hagoromo ikke? 

\says{D} Hvad er Hagoromo? 

\says{S} Altså det japanske kridt? 

\says{D} Kridt? 

\says{S} Ja du sælger vel Hagoromo kridt? 

\says{D} Nej? Jeg sælger coke? 

\scene{Lys ned}

\end{sketch}
\end{document}


